% --- Archon physics formalization source begin ---
% archon:physics
% archon:covers PhyXMiniProblems/problem_phyx_mini_0825.lean
% archon:source-report reports/phyx_mini/problem_phyx_mini_0825.source.json

\chapter{Physics problem phyx\_mini\_0825}
\label{ch:PhyXMiniProblems_problem_phyx_mini_0825}

\paragraph{Problem source.}
\#\# Physical scenario

A bowling ball of mass \textbackslash{}( M \textbackslash{}) rolls without slipping down a ramp that is inclined at an angle \textbackslash{}( \textbackslash{}beta \textbackslash{}) to the horizontal as shown in figure.

\#\# Question

What are the magnitude of the friction force on the ball? Treat the ball as a uniform solid sphere, ignoring the finger holes.

\#\# Answer choices

- A: A: \textbackslash{}( \textbackslash{}frac\{2\}\{9\} Mg \textbackslash{}sin \textbackslash{}beta \textbackslash{})

- B: B: \textbackslash{}( \textbackslash{}frac\{2\}\{7\} Mg \textbackslash{}sin \textbackslash{}beta \textbackslash{})

- C: C: \textbackslash{}( \textbackslash{}frac\{2\}\{5\} Mg \textbackslash{}cos \textbackslash{}beta \textbackslash{})

- D: D: \textbackslash{}( \textbackslash{}frac\{2\}\{3\} Mg \textbackslash{}cos \textbackslash{}beta \textbackslash{})

\#\# Figure caption (auxiliary; use the image as primary evidence)

The image depicts a bowling ball on an inclined plane. The scene shows the following elements:

- **Bowling Ball:** An orange sphere representing the bowling ball, located on the inclined plane. The ball has three finger holes.
  
- **Inclined Plane:** A gray surface on which the ball rests, inclined at an angle.
  
- **Angle β:** The angle of inclination between the plane and the horizontal surface, marked with a dashed line and arrow.

- **Label M:** Denotes the mass or a related concept associated with the ball.

This image likely illustrates concepts related to physics, such as motion on an incline or forces acting on the ball.

\#\# Dataset metadata

Dynamics; Physical Model Grounding Reasoning, Multi-Formula Reasoning

\paragraph{Recorded answer/context.}
B: B: \textbackslash{}( \textbackslash{}frac\{2\}\{7\} Mg \textbackslash{}sin \textbackslash{}beta \textbackslash{})

\paragraph{Figure/image path.}
/root/proposal\_for\_physic/science-mango/phyx\_mini\_run/phyx\_data/test\_image/825.png

\paragraph{Formalization target.}
create a compiling Lean file with sorry bodies at `PhyXMiniProblems/problem\_phyx\_mini\_0825.lean`. The Lean declarations must preserve the physical quantities, dimensions or dimensional roles, figure labels, governing-law hypotheses, and final relation expressed by this problem.
Use Mathlib/Physlib names found through LeanExplore where available. If a physics API is missing, introduce faithful local abstractions rather than scalar placeholder aliases.

\begin{theorem}[Physics formalization target]
\label{thm:physics:phyx_mini_0825:target}
\lean{PhyXMiniProblems.ProblemPhyXMini0825.frictionForceMagnitude_eq_twoSevenths_mass_gravity_sin_beta}
\uses{def:physics:phyx-mini-0825:phyxminiproblems-problemphyxmini0825-massinkilograms, def:physics:phyx-mini-0825:phyxminiproblems-problemphyxmini0825-accelerationmagnitudeinmeterspersecondsquared, def:physics:phyx-mini-0825:phyxminiproblems-problemphyxmini0825-signedforceinnewtons, def:physics:phyx-mini-0825:phyxminiproblems-problemphyxmini0825-rollingbowlingballsetup, def:physics:phyx-mini-0825:phyxminiproblems-problemphyxmini0825-matchesprimaryfigure, def:physics:phyx-mini-0825:phyxminiproblems-problemphyxmini0825-matchesbowlingballproblemdata, def:physics:phyx-mini-0825:phyxminiproblems-problemphyxmini0825-hasphysicalbowlingballparameters, def:physics:phyx-mini-0825:phyxminiproblems-problemphyxmini0825-satisfiesrollingspheredynamics}
The assigned autoformalize agent should translate this physics problem into Lean declarations in the covered file, with theorem and lemma proof bodies written as `by sorry`.
\end{theorem}
\begin{proof}
This is an autoformalization task, not a proof task. Produce statements that can later be proved by the physics prover without weakening the physical model.
\end{proof}

% --- Archon declaration topology begin ---
\section{Declaration topology}
\begin{definition}[Mass Quantity]
  \label{def:physics:phyx-mini-0825:phyxminiproblems-problemphyxmini0825-massquantity}
  \lean{PhyXMiniProblems.ProblemPhyXMini0825.MassQuantity}
  A nonnegative, unit-independent physical mass
\end{definition}
\begin{proof}
  This is the declared model definition of Mass Quantity.
\end{proof}
\begin{definition}[Length Quantity]
  \label{def:physics:phyx-mini-0825:phyxminiproblems-problemphyxmini0825-lengthquantity}
  \lean{PhyXMiniProblems.ProblemPhyXMini0825.LengthQuantity}
  A nonnegative physical length, used for the bowling-ball radius
\end{definition}
\begin{proof}
  This is the declared model definition of Length Quantity.
\end{proof}
\begin{definition}[acceleration Dimension]
  \label{def:physics:phyx-mini-0825:phyxminiproblems-problemphyxmini0825-accelerationdimension}
  \lean{PhyXMiniProblems.ProblemPhyXMini0825.accelerationDimension}
  The physical dimension of linear acceleration, L T⁻²
\end{definition}
\begin{proof}
  This is the declared model definition of acceleration Dimension.
\end{proof}
\begin{definition}[Acceleration Magnitude Quantity]
  \label{def:physics:phyx-mini-0825:phyxminiproblems-problemphyxmini0825-accelerationmagnitudequantity}
  \lean{PhyXMiniProblems.ProblemPhyXMini0825.AccelerationMagnitudeQuantity}
  \uses{def:physics:phyx-mini-0825:phyxminiproblems-problemphyxmini0825-accelerationdimension}
  A nonnegative gravitational-acceleration magnitude
\end{definition}
\begin{proof}
  This is the declared model definition of Acceleration Magnitude Quantity.
\end{proof}
\begin{definition}[Signed Acceleration Quantity]
  \label{def:physics:phyx-mini-0825:phyxminiproblems-problemphyxmini0825-signedaccelerationquantity}
  \lean{PhyXMiniProblems.ProblemPhyXMini0825.SignedAccelerationQuantity}
  \uses{def:physics:phyx-mini-0825:phyxminiproblems-problemphyxmini0825-accelerationdimension}
  A signed linear-acceleration component along the ramp
\end{definition}
\begin{proof}
  This is the declared model definition of Signed Acceleration Quantity.
\end{proof}
\begin{definition}[angular Acceleration Dimension]
  \label{def:physics:phyx-mini-0825:phyxminiproblems-problemphyxmini0825-angularaccelerationdimension}
  \lean{PhyXMiniProblems.ProblemPhyXMini0825.angularAccelerationDimension}
  The physical dimension of angular acceleration; radians are dimensionless
\end{definition}
\begin{proof}
  This is the declared model definition of angular Acceleration Dimension.
\end{proof}
\begin{definition}[Signed Angular Acceleration Quantity]
  \label{def:physics:phyx-mini-0825:phyxminiproblems-problemphyxmini0825-signedangularaccelerationquantity}
  \lean{PhyXMiniProblems.ProblemPhyXMini0825.SignedAngularAccelerationQuantity}
  \uses{def:physics:phyx-mini-0825:phyxminiproblems-problemphyxmini0825-angularaccelerationdimension}
  A signed angular-acceleration component about the rolling axis
\end{definition}
\begin{proof}
  This is the declared model definition of Signed Angular Acceleration Quantity.
\end{proof}
\begin{definition}[force Dimension]
  \label{def:physics:phyx-mini-0825:phyxminiproblems-problemphyxmini0825-forcedimension}
  \lean{PhyXMiniProblems.ProblemPhyXMini0825.forceDimension}
  The physical dimension of force, M L T⁻²
\end{definition}
\begin{proof}
  This is the declared model definition of force Dimension.
\end{proof}
\begin{definition}[Signed Force Quantity]
  \label{def:physics:phyx-mini-0825:phyxminiproblems-problemphyxmini0825-signedforcequantity}
  \lean{PhyXMiniProblems.ProblemPhyXMini0825.SignedForceQuantity}
  \uses{def:physics:phyx-mini-0825:phyxminiproblems-problemphyxmini0825-forcedimension}
  A signed force component along the ramp
\end{definition}
\begin{proof}
  This is the declared model definition of Signed Force Quantity.
\end{proof}
\begin{definition}[torque Dimension]
  \label{def:physics:phyx-mini-0825:phyxminiproblems-problemphyxmini0825-torquedimension}
  \lean{PhyXMiniProblems.ProblemPhyXMini0825.torqueDimension}
  \uses{def:physics:phyx-mini-0825:phyxminiproblems-problemphyxmini0825-forcedimension}
  The physical dimension of torque, M L² T⁻²
\end{definition}
\begin{proof}
  This is the declared model definition of torque Dimension.
\end{proof}
\begin{definition}[Signed Torque Quantity]
  \label{def:physics:phyx-mini-0825:phyxminiproblems-problemphyxmini0825-signedtorquequantity}
  \lean{PhyXMiniProblems.ProblemPhyXMini0825.SignedTorqueQuantity}
  \uses{def:physics:phyx-mini-0825:phyxminiproblems-problemphyxmini0825-torquedimension}
  A signed torque component about the rolling axis
\end{definition}
\begin{proof}
  This is the declared model definition of Signed Torque Quantity.
\end{proof}
\begin{definition}[moment Of Inertia Dimension]
  \label{def:physics:phyx-mini-0825:phyxminiproblems-problemphyxmini0825-momentofinertiadimension}
  \lean{PhyXMiniProblems.ProblemPhyXMini0825.momentOfInertiaDimension}
  The physical dimension of an axial moment of inertia, M L²
\end{definition}
\begin{proof}
  This is the declared model definition of moment Of Inertia Dimension.
\end{proof}
\begin{definition}[Moment Of Inertia Quantity]
  \label{def:physics:phyx-mini-0825:phyxminiproblems-problemphyxmini0825-momentofinertiaquantity}
  \lean{PhyXMiniProblems.ProblemPhyXMini0825.MomentOfInertiaQuantity}
  \uses{def:physics:phyx-mini-0825:phyxminiproblems-problemphyxmini0825-momentofinertiadimension}
  A nonnegative scalar moment of inertia about the rolling axis
\end{definition}
\begin{proof}
  This is the declared model definition of Moment Of Inertia Quantity.
\end{proof}
\begin{definition}[nonnegative SIReadout]
  \label{def:physics:phyx-mini-0825:phyxminiproblems-problemphyxmini0825-nonnegativesireadout}
  \lean{PhyXMiniProblems.ProblemPhyXMini0825.nonnegativeSIReadout}
  Read a nonnegative dimensionful quantity in coherent SI units
\end{definition}
\begin{proof}
  This is the declared model definition of nonnegative SIReadout.
\end{proof}
\begin{definition}[signed SIReadout]
  \label{def:physics:phyx-mini-0825:phyxminiproblems-problemphyxmini0825-signedsireadout}
  \lean{PhyXMiniProblems.ProblemPhyXMini0825.signedSIReadout}
  Read a signed dimensionful scalar component in coherent SI units
\end{definition}
\begin{proof}
  This is the declared model definition of signed SIReadout.
\end{proof}
\begin{definition}[mass In Kilograms]
  \label{def:physics:phyx-mini-0825:phyxminiproblems-problemphyxmini0825-massinkilograms}
  \lean{PhyXMiniProblems.ProblemPhyXMini0825.massInKilograms}
  \uses{def:physics:phyx-mini-0825:phyxminiproblems-problemphyxmini0825-massquantity, def:physics:phyx-mini-0825:phyxminiproblems-problemphyxmini0825-nonnegativesireadout}
  Kilogram readout of the ball's mass
\end{definition}
\begin{proof}
  This is the declared model definition of mass In Kilograms.
\end{proof}
\begin{definition}[length In Meters]
  \label{def:physics:phyx-mini-0825:phyxminiproblems-problemphyxmini0825-lengthinmeters}
  \lean{PhyXMiniProblems.ProblemPhyXMini0825.lengthInMeters}
  \uses{def:physics:phyx-mini-0825:phyxminiproblems-problemphyxmini0825-lengthquantity, def:physics:phyx-mini-0825:phyxminiproblems-problemphyxmini0825-nonnegativesireadout}
  Metre readout of the ball's radius
\end{definition}
\begin{proof}
  This is the declared model definition of length In Meters.
\end{proof}
\begin{definition}[acceleration Magnitude In Meters Per Second Squared]
  \label{def:physics:phyx-mini-0825:phyxminiproblems-problemphyxmini0825-accelerationmagnitudeinmeterspersecondsquared}
  \lean{PhyXMiniProblems.ProblemPhyXMini0825.accelerationMagnitudeInMetersPerSecondSquared}
  \uses{def:physics:phyx-mini-0825:phyxminiproblems-problemphyxmini0825-accelerationmagnitudequantity, def:physics:phyx-mini-0825:phyxminiproblems-problemphyxmini0825-nonnegativesireadout}
  Metre-per-second-squared readout of a nonnegative acceleration
\end{definition}
\begin{proof}
  This is the declared model definition of acceleration Magnitude In Meters Per Second Squared.
\end{proof}
\begin{definition}[signed Acceleration In Meters Per Second Squared]
  \label{def:physics:phyx-mini-0825:phyxminiproblems-problemphyxmini0825-signedaccelerationinmeterspersecondsquared}
  \lean{PhyXMiniProblems.ProblemPhyXMini0825.signedAccelerationInMetersPerSecondSquared}
  \uses{def:physics:phyx-mini-0825:phyxminiproblems-problemphyxmini0825-signedaccelerationquantity, def:physics:phyx-mini-0825:phyxminiproblems-problemphyxmini0825-signedsireadout}
  Metre-per-second-squared readout of a signed rampwise acceleration
\end{definition}
\begin{proof}
  This is the declared model definition of signed Acceleration In Meters Per Second Squared.
\end{proof}
\begin{definition}[angular Acceleration In Radians Per Second Squared]
  \label{def:physics:phyx-mini-0825:phyxminiproblems-problemphyxmini0825-angularaccelerationinradianspersecondsquared}
  \lean{PhyXMiniProblems.ProblemPhyXMini0825.angularAccelerationInRadiansPerSecondSquared}
  \uses{def:physics:phyx-mini-0825:phyxminiproblems-problemphyxmini0825-signedangularaccelerationquantity, def:physics:phyx-mini-0825:phyxminiproblems-problemphyxmini0825-signedsireadout}
  Radian-per-second-squared readout of a signed angular acceleration
\end{definition}
\begin{proof}
  This is the declared model definition of angular Acceleration In Radians Per Second Squared.
\end{proof}
\begin{definition}[signed Force In Newtons]
  \label{def:physics:phyx-mini-0825:phyxminiproblems-problemphyxmini0825-signedforceinnewtons}
  \lean{PhyXMiniProblems.ProblemPhyXMini0825.signedForceInNewtons}
  \uses{def:physics:phyx-mini-0825:phyxminiproblems-problemphyxmini0825-signedforcequantity, def:physics:phyx-mini-0825:phyxminiproblems-problemphyxmini0825-signedsireadout}
  Newton readout of a signed rampwise force component
\end{definition}
\begin{proof}
  This is the declared model definition of signed Force In Newtons.
\end{proof}
\begin{definition}[signed Torque In Newton Meters]
  \label{def:physics:phyx-mini-0825:phyxminiproblems-problemphyxmini0825-signedtorqueinnewtonmeters}
  \lean{PhyXMiniProblems.ProblemPhyXMini0825.signedTorqueInNewtonMeters}
  \uses{def:physics:phyx-mini-0825:phyxminiproblems-problemphyxmini0825-signedtorquequantity, def:physics:phyx-mini-0825:phyxminiproblems-problemphyxmini0825-signedsireadout}
  Newton-metre readout of a signed torque component
\end{definition}
\begin{proof}
  This is the declared model definition of signed Torque In Newton Meters.
\end{proof}
\begin{definition}[moment Of Inertia In Kilogram Meters Squared]
  \label{def:physics:phyx-mini-0825:phyxminiproblems-problemphyxmini0825-momentofinertiainkilogrammeterssquared}
  \lean{PhyXMiniProblems.ProblemPhyXMini0825.momentOfInertiaInKilogramMetersSquared}
  \uses{def:physics:phyx-mini-0825:phyxminiproblems-problemphyxmini0825-momentofinertiaquantity, def:physics:phyx-mini-0825:phyxminiproblems-problemphyxmini0825-nonnegativesireadout}
  Kilogram-metre-squared readout of a scalar axial moment of inertia
\end{definition}
\begin{proof}
  This is the declared model definition of moment Of Inertia In Kilogram Meters Squared.
\end{proof}
\begin{definition}[Figure Part]
  \label{def:physics:phyx-mini-0825:phyxminiproblems-problemphyxmini0825-figurepart}
  \lean{PhyXMiniProblems.ProblemPhyXMini0825.FigurePart}
  Visible geometric or physical parts of the supplied bitmap
\end{definition}
\begin{proof}
  This is the declared model definition of Figure Part.
\end{proof}
\begin{definition}[Figure Label]
  \label{def:physics:phyx-mini-0825:phyxminiproblems-problemphyxmini0825-figurelabel}
  \lean{PhyXMiniProblems.ProblemPhyXMini0825.FigureLabel}
  Literal symbolic labels printed in the supplied bitmap
\end{definition}
\begin{proof}
  This is the declared model definition of Figure Label.
\end{proof}
\begin{definition}[Depicted Ball Placement]
  \label{def:physics:phyx-mini-0825:phyxminiproblems-problemphyxmini0825-depictedballplacement}
  \lean{PhyXMiniProblems.ProblemPhyXMini0825.DepictedBallPlacement}
  The qualitative placement of the depicted ball
\end{definition}
\begin{proof}
  This is the declared model definition of Depicted Ball Placement.
\end{proof}
\begin{definition}[Bowling Ball Incline Figure]
  \label{def:physics:phyx-mini-0825:phyxminiproblems-problemphyxmini0825-bowlingballinclinefigure}
  \lean{PhyXMiniProblems.ProblemPhyXMini0825.BowlingBallInclineFigure}
  \uses{def:physics:phyx-mini-0825:phyxminiproblems-problemphyxmini0825-figurepart, def:physics:phyx-mini-0825:phyxminiproblems-problemphyxmini0825-figurelabel, def:physics:phyx-mini-0825:phyxminiproblems-problemphyxmini0825-depictedballplacement}
  Qualitative and symbolic information read directly from image 825.png
\end{definition}
\begin{proof}
  This is the declared model definition of Bowling Ball Incline Figure.
\end{proof}
\begin{definition}[Ball Model]
  \label{def:physics:phyx-mini-0825:phyxminiproblems-problemphyxmini0825-ballmodel}
  \lean{PhyXMiniProblems.ProblemPhyXMini0825.BallModel}
  The rigid-body idealization stipulated in the problem text
\end{definition}
\begin{proof}
  This is the declared model definition of Ball Model.
\end{proof}
\begin{definition}[Contact Regime]
  \label{def:physics:phyx-mini-0825:phyxminiproblems-problemphyxmini0825-contactregime}
  \lean{PhyXMiniProblems.ProblemPhyXMini0825.ContactRegime}
  Tangential contact behavior between the ball and ramp
\end{definition}
\begin{proof}
  This is the declared model definition of Contact Regime.
\end{proof}
\begin{definition}[Tangential Positive Direction]
  \label{def:physics:phyx-mini-0825:phyxminiproblems-problemphyxmini0825-tangentialpositivedirection}
  \lean{PhyXMiniProblems.ProblemPhyXMini0825.TangentialPositiveDirection}
  Orientation used for signed tangential scalar components
\end{definition}
\begin{proof}
  This is the declared model definition of Tangential Positive Direction.
\end{proof}
\begin{definition}[Angular Positive Direction]
  \label{def:physics:phyx-mini-0825:phyxminiproblems-problemphyxmini0825-angularpositivedirection}
  \lean{PhyXMiniProblems.ProblemPhyXMini0825.AngularPositiveDirection}
  Orientation used for signed angular scalar components
\end{definition}
\begin{proof}
  This is the declared model definition of Angular Positive Direction.
\end{proof}
\begin{definition}[Rolling Bowling Ball Setup]
  \label{def:physics:phyx-mini-0825:phyxminiproblems-problemphyxmini0825-rollingbowlingballsetup}
  \lean{PhyXMiniProblems.ProblemPhyXMini0825.RollingBowlingBallSetup}
  \uses{def:physics:phyx-mini-0825:phyxminiproblems-problemphyxmini0825-massquantity, def:physics:phyx-mini-0825:phyxminiproblems-problemphyxmini0825-lengthquantity, def:physics:phyx-mini-0825:phyxminiproblems-problemphyxmini0825-accelerationmagnitudequantity, def:physics:phyx-mini-0825:phyxminiproblems-problemphyxmini0825-signedaccelerationquantity, def:physics:phyx-mini-0825:phyxminiproblems-problemphyxmini0825-signedangularaccelerationquantity, def:physics:phyx-mini-0825:phyxminiproblems-problemphyxmini0825-signedforcequantity, def:physics:phyx-mini-0825:phyxminiproblems-problemphyxmini0825-signedtorquequantity, def:physics:phyx-mini-0825:phyxminiproblems-problemphyxmini0825-momentofinertiaquantity, def:physics:phyx-mini-0825:phyxminiproblems-problemphyxmini0825-bowlingballinclinefigure, def:physics:phyx-mini-0825:phyxminiproblems-problemphyxmini0825-ballmodel, def:physics:phyx-mini-0825:phyxminiproblems-problemphyxmini0825-contactregime, def:physics:phyx-mini-0825:phyxminiproblems-problemphyxmini0825-tangentialpositivedirection, def:physics:phyx-mini-0825:phyxminiproblems-problemphyxmini0825-angularpositivedirection}
  Every unknown observable is an independent physical quantity. In particular, the friction component is not defined from an answer choice. rigidBodySI uses Physlib's SI-coordinate rigid-body interface, while the corresponding mass, radius, force, acceleration, torque, and scalar inertia remain unit-independent quantities
\end{definition}
\begin{proof}
  This is the declared model definition of Rolling Bowling Ball Setup.
\end{proof}
\begin{definition}[Matches Primary Figure]
  \label{def:physics:phyx-mini-0825:phyxminiproblems-problemphyxmini0825-matchesprimaryfigure}
  \lean{PhyXMiniProblems.ProblemPhyXMini0825.MatchesPrimaryFigure}
  \uses{def:physics:phyx-mini-0825:phyxminiproblems-problemphyxmini0825-figurepart, def:physics:phyx-mini-0825:phyxminiproblems-problemphyxmini0825-figurelabel, def:physics:phyx-mini-0825:phyxminiproblems-problemphyxmini0825-rollingbowlingballsetup}
  Labels, geometry, and placement transcribed from the primary bitmap
\end{definition}
\begin{proof}
  This is the declared model definition of Matches Primary Figure.
\end{proof}
\begin{definition}[Matches Bowling Ball Problem Data]
  \label{def:physics:phyx-mini-0825:phyxminiproblems-problemphyxmini0825-matchesbowlingballproblemdata}
  \lean{PhyXMiniProblems.ProblemPhyXMini0825.MatchesBowlingBallProblemData}
  \uses{def:physics:phyx-mini-0825:phyxminiproblems-problemphyxmini0825-rollingbowlingballsetup}
  The text identifies M as the ball's mass, specifies rolling without slip, and replaces the visibly perforated bowling ball by a uniform solid sphere. The equality to RigidBody.solidSphere is model identification, not an assumption about the requested friction force
\end{definition}
\begin{proof}
  This is the declared model definition of Matches Bowling Ball Problem Data.
\end{proof}
\begin{definition}[Has Physical Bowling Ball Parameters]
  \label{def:physics:phyx-mini-0825:phyxminiproblems-problemphyxmini0825-hasphysicalbowlingballparameters}
  \lean{PhyXMiniProblems.ProblemPhyXMini0825.HasPhysicalBowlingBallParameters}
  \uses{def:physics:phyx-mini-0825:phyxminiproblems-problemphyxmini0825-massinkilograms, def:physics:phyx-mini-0825:phyxminiproblems-problemphyxmini0825-lengthinmeters, def:physics:phyx-mini-0825:phyxminiproblems-problemphyxmini0825-accelerationmagnitudeinmeterspersecondsquared, def:physics:phyx-mini-0825:phyxminiproblems-problemphyxmini0825-rollingbowlingballsetup}
  Positivity and angle bounds selecting the depicted physical branch
\end{definition}
\begin{proof}
  This is the declared model definition of Has Physical Bowling Ball Parameters.
\end{proof}
\begin{definition}[Satisfies Rolling Sphere Dynamics]
  \label{def:physics:phyx-mini-0825:phyxminiproblems-problemphyxmini0825-satisfiesrollingspheredynamics}
  \lean{PhyXMiniProblems.ProblemPhyXMini0825.SatisfiesRollingSphereDynamics}
  \uses{def:physics:phyx-mini-0825:phyxminiproblems-problemphyxmini0825-massinkilograms, def:physics:phyx-mini-0825:phyxminiproblems-problemphyxmini0825-lengthinmeters, def:physics:phyx-mini-0825:phyxminiproblems-problemphyxmini0825-accelerationmagnitudeinmeterspersecondsquared, def:physics:phyx-mini-0825:phyxminiproblems-problemphyxmini0825-signedaccelerationinmeterspersecondsquared, def:physics:phyx-mini-0825:phyxminiproblems-problemphyxmini0825-angularaccelerationinradianspersecondsquared, def:physics:phyx-mini-0825:phyxminiproblems-problemphyxmini0825-signedforceinnewtons, def:physics:phyx-mini-0825:phyxminiproblems-problemphyxmini0825-signedtorqueinnewtonmeters, def:physics:phyx-mini-0825:phyxminiproblems-problemphyxmini0825-momentofinertiainkilogrammeterssquared, def:physics:phyx-mini-0825:phyxminiproblems-problemphyxmini0825-rollingbowlingballsetup}
  The signed force and acceleration components are positive down the ramp, and the signed angular quantities are positive in the corresponding rolling-down direction. Thus the torque law contains a minus sign: a positive down-ramp contact force produces torque opposite the rolling-down orientation. No field below contains the solved friction value or assumes that friction is uphill.
\end{definition}
\begin{proof}
  This is the declared model definition of Satisfies Rolling Sphere Dynamics.
\end{proof}
\begin{lemma}[axial Moment Of Inertia eq two Fifths mass mul radius sq]
  \label{lem:physics:phyx-mini-0825:phyxminiproblems-problemphyxmini0825-axialmomentofinertia-eq-twofifths-mass-mul-radius-sq}
  \lean{PhyXMiniProblems.ProblemPhyXMini0825.axialMomentOfInertia_eq_twoFifths_mass_mul_radius_sq}
  \uses{def:physics:phyx-mini-0825:phyxminiproblems-problemphyxmini0825-massinkilograms, def:physics:phyx-mini-0825:phyxminiproblems-problemphyxmini0825-lengthinmeters, def:physics:phyx-mini-0825:phyxminiproblems-problemphyxmini0825-momentofinertiainkilogrammeterssquared, def:physics:phyx-mini-0825:phyxminiproblems-problemphyxmini0825-rollingbowlingballsetup, def:physics:phyx-mini-0825:phyxminiproblems-problemphyxmini0825-matchesbowlingballproblemdata, def:physics:phyx-mini-0825:phyxminiproblems-problemphyxmini0825-hasphysicalbowlingballparameters, def:physics:phyx-mini-0825:phyxminiproblems-problemphyxmini0825-satisfiesrollingspheredynamics}
  Physlib's uniform-solid-sphere tensor formula gives the familiar scalar axial moment of inertia I = (2/5) M R² for the modeled bowling ball
\end{lemma}
\begin{proof}
  The conclusion follows from the governing relations and figure readouts named in the statement of axial Moment Of Inertia eq two Fifths mass mul radius sq.
\end{proof}
\begin{definition}[Answer Choice]
  \label{def:physics:phyx-mini-0825:phyxminiproblems-problemphyxmini0825-answerchoice}
  \lean{PhyXMiniProblems.ProblemPhyXMini0825.AnswerChoice}
  Labels of the four displayed multiple-choice answers
\end{definition}
\begin{proof}
  This is the declared model definition of Answer Choice.
\end{proof}
\begin{definition}[displayed Magnitude In Newtons]
  \label{def:physics:phyx-mini-0825:phyxminiproblems-problemphyxmini0825-answerchoice-displayedmagnitudeinnewtons}
  \lean{PhyXMiniProblems.ProblemPhyXMini0825.AnswerChoice.displayedMagnitudeInNewtons}
  \uses{def:physics:phyx-mini-0825:phyxminiproblems-problemphyxmini0825-massinkilograms, def:physics:phyx-mini-0825:phyxminiproblems-problemphyxmini0825-accelerationmagnitudeinmeterspersecondsquared, def:physics:phyx-mini-0825:phyxminiproblems-problemphyxmini0825-rollingbowlingballsetup, def:physics:phyx-mini-0825:phyxminiproblems-problemphyxmini0825-answerchoice}
  The force magnitudes printed in the four answer choices. This records the question's answer metadata only; it does not identify any choice as correct
\end{definition}
\begin{proof}
  This is the declared model definition of displayed Magnitude In Newtons.
\end{proof}
% --- Archon declaration topology end ---
% --- Archon physics formalization source end ---
